\section{Introduction}
\label{sec:intro}
Audio temporal grounding identifies the intervals in an audio recording that match a natural-language query. Text-to-audio grounding established phrase-level temporal correspondence~\cite{xu2021grounding}, and audio moment retrieval extended query-based localization to untrimmed recordings~\cite{munakata2025amr}. SpotSound uses timestamp-interleaved audio representations and training on absent events to strengthen audio-language grounding~\cite{sun2026spotsound}. These capabilities motivate a complementary question: once a strong model has generated an interval, when should its boundaries be revised?

An acoustic transition is evidence about a boundary, but not necessarily evidence that moving the current boundary will improve localization. Shrinking a predicted interval can remove background or truncate a correctly identified event. A fixed displacement also changes a larger fraction of a short event than of a long one. A useful refinement rule must therefore address both the value of a candidate edit and whether the available evidence warrants replacing the original answer.

Refinement itself is established in temporal grounding. TimeRefine predicts successive offsets from coarse video intervals~\cite{wang2024timerefine}; TimePLE includes duration-aware bounded correction within an interval representation~\cite{zeng2026timeple}. We focus instead on \emph{incumbent-relative edit utility}: a candidate is evaluated against the existing prediction, and retaining that prediction is an explicit alternative. This separates a plausible boundary proposal from an edit worth applying, without claiming that boundary refinement or adaptive radii are new in isolation.

We implement this formulation as a selective boundary-refinement extension to \NOVA{} (Necessity-Oriented Verification over Audio). Frozen acoustic features support a lightweight supervised utility model. A constrained dynamic program combines legal edits, and a bootstrap-ensemble criterion selects between the proposed update and the identity action. Duration-dependent radii restrict displacement for short events. The criterion is an empirical conservative decision rule, not a distribution-free safety guarantee.

On SpotSound-Bench and Clotho-Moment, the system reaches 59.39\% and 86.86\% mIoU, respectively. Relative to the best entries in the cited SpotSound main table, the differences are 1.49 and 1.26 points. We distinguish these system-level results from component attribution and separate primary comparisons, protocol-qualified references, and development evidence. Our methodological focus is the combination of relative utility learning, structured edit selection, and explicit retention of an existing answer.
